\section{Discussion and Conclusion} \label{conclusion}

This paper asks two questions about how the prospect of punishment interacts with the use of algorithms as decision aids. Can decision-makers shift responsibility for bad outcomes onto algorithms by delegating their decisions? And does the anticipation of being held responsible motivate delegation in the first place? We address both questions in a preregistered online experiment in which Player~A makes a payoff-relevant prediction for Player~B and may delegate the prediction to an algorithm calibrated to perform as well as Player~A herself. A between-subject manipulation switches off punishment by Player~B in the \textit{No-Punishment} condition; the \textit{Punishment} preserves the natural setting in which the prospect of punishment is on the table.

Two findings emerge. First, and most striking, the \textit{prospect} of punishment \textit{reduces}, rather than increases, the propensity to delegate to the algorithm: $42.5\%$ of Player~As delegate when punishment is possible, compared to $59.3\%$ when it is not. This direction is the opposite of what the responsibility-avoidance literature would predict, and we view the reversal as the paper's main empirical contribution---noting that, because the result reverses the preregistered direction, the analyses we use to interpret it (Section~\ref{sec:mechanism}) are exploratory rather than confirmatory. Second, conditional on the outcome for Player~B, we cannot reject the hypothesis that delegated and self-made decisions are punished equally. The point estimate runs in the predicted direction (lower punishment under delegation in the bad-outcome cell) but is small relative to the test's confidence interval at the achieved sample size; we therefore frame Result~2 as a failure to detect insulation rather than as positive evidence that delegation cannot insulate. Player~A does not appear to be able to use this particular algorithm as a meaningful blame shield in the bounded sense the test can assess.

\paragraph{Mechanism.} Why does punishment reduce, rather than increase, delegation? The data are inconsistent with a simple anticipation mechanism: Player~As do not on average expect harsher punishment for delegated decisions, and they correctly anticipate the absence of differential punishment that we observe in Player~B's actual choices. The data are likewise inconsistent with the simple form of an effort story in which the prospect of punishment motivates Player~A to outperform the algorithm: among non-delegators, those in the \textit{No-Punishment} condition (where punishment is impossible) outperform those in the \textit{Punishment} condition, the reverse of what such a story implies. The reading most consistent with the data is one of \textit{process ownership}: Player~As who actively choose not to delegate take responsibility for the outcome, exert correspondingly higher effort, and that share is higher when punishment is on the table. The supporting evidence is exploratory rather than confirmatory, however, and rests on a single performance-by-delegation heterogeneity pattern in a 60-subject sub-sample (Column (3) of Table~\ref{tab:del_decision_determinants}); we cannot, with this design, separate process ownership cleanly from an experimenter-demand reading in which the punishment cue itself nudges Player~A toward direct engagement. We therefore present process ownership as the most consistent reading of the data we have, not as a confirmed mechanism. Disentangling it from demand is a clean direction for follow-up work, sketched below.

This refinement is consistent with, but goes beyond, the standard reading of the responsibility-avoidance literature. The classic finding---that delegation reduces punishment of the principal and that this prospect motivates delegation---was established in dictator-game settings where the underlying choice is between a known fair and a known unfair allocation, the intermediary is human, and the principal can be held accountable specifically for choosing a particular intermediary~\citep{bartling_shifting_2012,coffman_intermediation_2011,oexl_shifting_2013}. When the intermediary is an algorithm calibrated to mimic the principal and the underlying decision involves genuine uncertainty, none of the standard mechanisms through which delegation can shift responsibility---signal value of the delegation choice, the human-human punishment channel, the certainty of the unfair outcome---are available. What remains is a baseline outcome-based punishment, supplemented by an asymmetric desire to be the proximate cause of good outcomes when the moral stakes are salient.

\paragraph{Implications, with explicit bounds.} Our findings speak to the design of accountability regimes governing algorithmic decision-making. The European Union's AI~Act~\citep{european_union_ai_act_2024} and analogous regulatory frameworks elsewhere place increasing emphasis on assigning responsibility to human users of high-risk algorithmic systems. The implicit theory of behaviour underlying these provisions is that holding humans accountable for algorithm-mediated decisions \textit{disciplines} their use of those algorithms. Our results offer a behavioural complement to this view: when punishment is on the table, decision-makers in our setting do not simply delegate-and-blame the algorithm; they choose to take direct responsibility for the decision more often. From the perspective of efficiency, this need not be welfare-increasing: in domains where algorithms genuinely outperform humans, holding humans liable for outcomes may push decision-makers away from the more accurate technology. The trade-off between encouraging algorithmic adoption and assigning meaningful responsibility is, on this reading, sharper than the existing literature has suggested.

We note explicitly the bounds of this reading. Our evidence is a single online experiment with $322$ UK Prolific participants on a stylised weight-prediction task, with stakes of $\pounds 5$ for Player~B and a punishment range of $\pounds 0$--$\pounds 2$. The high-stakes algorithmic domains that motivate the regulatory debate---medical diagnosis, judicial risk assessment, hiring, credit allocation---differ from our setting in stake magnitude, professional context, expert training, the institutional structure of accountability, and the kind of population at the decision margin. We therefore view the result as identifying a behavioural channel that exists in principle and is detectable in our setting, not as a calibration of how big a lever accountability is in any specific applied domain. Mapping the magnitude of the effect against external benchmarks---e.g., physician adoption of diagnostic algorithms under malpractice exposure, or loan-officer deviation from credit-scoring algorithms under fair-lending audit---would require either domain-specific replication or instruments to which we do not have access. We flag this gap explicitly so that readers can calibrate the policy implication appropriately.

\paragraph{Limitations.} Several caveats are in order. The belief measure we elicit is unincentivised and noisy; it conveys ordinal patterns reliably but cannot support fine-grained quantitative claims about Player~A's anticipations. Two of the four belief scenarios are hypothetical conditional on the delegation decision, and all four are hypothetical in the \textit{No-Punishment} condition; this further limits the inference, and we plan to report the H2 patterns separately for the payoff-relevant and the hypothetical cells once the analysis pipeline is in place. We use the strategy method for Player~B's punishment decision, which is standard in this literature but may yield different absolute levels of punishment than a direct elicitation. Player~A also makes the eleventh prediction whether or not she ultimately decides to delegate, which equates the time cost of the two options but may signal her commitment differently than a design in which the prediction is skipped under delegation. We did not include a separate manipulation-comprehension check verifying that Player~A internalised the punishment-on/off distinction; we relied on the prominent screen-level framing and on attention-check pass rates as indirect evidence, but a future replication should add this check. The absolute stakes are modest (£5 high payoff, £0--2 punishment range): we cannot speak to how the effect scales with stakes, and a stake-magnitude robustness check is the first item in the future-research list below. Our sample is drawn from UK Prolific participants, which limits external validity along the dimensions in which Prolific samples differ from broader populations and from professional decision-makers in the high-stakes algorithmic domains we discuss above.

A natural concern is whether the punishment possibility cues Player~A to behave more conservatively for reasons unrelated to the substantive question---an experimenter-demand effect. We cannot rule this concern out from the data. The direction of any such effect is also ambiguous: a decision-maker who is reminded of being judged could plausibly delegate \textit{more} (to shed responsibility) or \textit{less} (to demonstrate effort and commitment); the literature does not unambiguously favour one over the other. We note two patterns weakly inconsistent with a strong demand reading---behaviour does not track elicited beliefs as tightly as a pure-demand story would predict, and the sharp performance heterogeneity in the \textit{Punishment} condition (better performers delegating less) is harder to motivate from demand alone---but we treat these as suggestive rather than dispositive. A direct test in the spirit of \citet{de_quidt_measuring_2018} would be the natural next step.

\paragraph{Directions for future research.} Several extensions would strengthen the inference. A direct-method companion would test the robustness of Result~2 against the strategy-method elicitation, and a meaningfully larger sample would let us put a tight confidence interval on the small but persistent point estimate of insulation under delegation. A condition in which Player~A is told explicitly that Player~B is unaware of the equal-performance calibration would test whether the relative-performance signal is the binding margin distinguishing our results from \citet{feier_hiding_2022}. Conditions that experimentally vary the algorithm's performance relative to Player~A (above, equal, below) would map the generality of the process-ownership reading. A demand-effect probe along the lines of \citet{de_quidt_measuring_2018} would put a credible upper bound on the contribution of M4. A stake-magnitude robustness check, varying the £5/£1 high/low payoffs and the £0--£2 punishment range, would speak to whether the effect scales with stakes in the direction that policy-relevant settings would imply. Finally, replication in a different country, in a different decision domain (e.g., a medical-triage or hiring vignette), or with a sample of professional decision-makers rather than online crowd-workers would speak directly to the external-validity bounds we describe above.

Taken together, our results suggest a more textured story than the one that has emerged from the existing literature on responsibility avoidance and algorithm aversion. Within the bounds of the setting we study, the capacity to use algorithms as a tool for shifting responsibility appears limited, and the prospect of being held responsible---rather than encouraging delegation as a hedge---appears to push decision-makers towards direct engagement with the consequential choice. Whether the same channel operates in the high-stakes algorithmic domains that animate the regulatory debate is an empirical question that this study can frame but cannot, on its own, answer.
